% Attestations and Compliance Statements (MAIN POINT 2). Movement 1 sets the
% duties and gives a model attestation block; Movement 2 maps each clause to the
% machine-checkable VVUQ-02 evidence in one compliance-items table.

This bill requires two instruments before any robot-patient interaction code is
generated or executed: an artificial-intelligence compliance statement and a
signed attestation by a responsible officer of the sponsor. Movement 1 sets out
their content; Movement 2 shows that each attested clause is already
machine-checkable against the recorded second-study evidence, so the attestation
certifies facts a reviewer can independently confirm rather than promises a
reviewer must trust.

\subsection{The Attestation and Compliance-Statement Duties}

The compliance statement is modelled on the Texas instrument: it states a summary
of the artificial-intelligence system's function, a decision-making logic tree, a
description of the training data sets, and a formal attestation that the system
complies with bias-reduction and evidence-based clinical criteria. The signed
attestation then certifies, under the penalties specified in
Section~\ref{sec:implementation_enforcement}, each of the following: that the
VVUQ verification process cleared before code generation, which is the
verification-before-generation certification at the core of this bill; that each
gate met its threshold bound to a named external standard; that the site passed
its Physical AI Standard Level gate and each robot met the Unification Standard
Level minimum for the procedure type; that the algorithm and its training data
minimize the risk of bias on protected characteristics and adhere to
evidence-based clinical guidelines, consistent with the New York certification
\cite{ny-a9149} and the federal nondiscrimination guarantees
\cite{usc-titlevi, usc-ada}; and that a licensed human retained the
medical-necessity decision, consistent with the California physician-decision
rule \cite{ca-sb1120} and the Florida human-review requirement \cite{fl-hb527}.
The integrity of the attested records is tied to the HIPAA and electronic-records
rules \cite{cfr-hipaa, cfr-part11} and to the risk-management standard
\cite{iso14971}. A model attestation block follows; the full bill may refine its
wording, but the certified facts are fixed.

\begin{quote}
\small\itshape
I, the undersigned responsible officer of the sponsor, certify under the
penalties of Section~\ref{sec:implementation_enforcement} that, for each
robot-patient interaction algorithm covered by this submission: (1) the automated
VVUQ verification process cleared before any covered code was generated or
executed; (2) each VVUQ gate met its threshold bound to a named external
standard; (3) the trial site passed its Physical AI Standard Level gate and each
robot met the Unification Standard Level minimum for its procedure type; (4) the
algorithm and its training data minimize the risk of bias on protected
characteristics and adhere to evidence-based clinical guidelines; and (5) a
licensed human professional retained the medical-necessity decision. The
supporting records are referenced, not attached, and are reproducible from the
stated seed.\par
\medskip
\upshape\rule{6cm}{0.4pt}\quad\rule{3.2cm}{0.4pt}\\
\small Responsible officer (signature)\hspace{1.4cm} Date
\end{quote}

\subsection{Demonstration from the Recorded Compliance Evidence}

Each certified clause maps to a recorded artifact that evidences it and to the
standard it answers to, as set out in Table~\ref{tab:attest-items}. The
verification-cleared-first clause is evidenced by the gate specification, which
fixes the verification fraction at 1.0 in the spirit of the ASME V\&V 40
credibility process. The gate-bound-to-a-standard clause is evidenced by the
standards map, which records the governing standard and enforced limit for each
of the ten gates. The continuous-integration clause is evidenced by the project
compliance checklist, which runs a lint pass, the 172-test suite, schema
validation, a YAML lint, and pre-commit hooks, and exercises the lint-and-format
matrix across Python 3.10, 3.11, and 3.12, tracking the medical-device software
life-cycle posture \cite{iec62304}. The readiness clause is evidenced by the
national-platform standards section, which records the site Physical AI Standard
Level pass or fail by dimension and the robot Unification Standard Level score by
dimension \cite{kawchak2026national}. Because these checks are automated and
reproducible, the attestation is an assertion about verifiable machine output,
consistent with the agency credibility framework \cite{fda2023credibility,
kawchak2026vvuq02}.

\begin{table}[ht]
\centering
\begin{tabularx}{\textwidth}{L{4.4cm} L{4.0cm} Y}
\toprule
\textbf{Attestation clause} & \textbf{Evidencing artifact (abbreviated)} & \textbf{Answers to} \\
\midrule
\multicolumn{3}{l}{\textit{Parent: papers/VVUQ-02/codegen/docs/}} \\
Verification cleared first & vvuq\_gate\_spec.md & ASME V\&V 40 § 8; verification fraction 1.0 \\
Gate bound to a standard & standards\_map.md & IEC 80601-2-77; ISO 14971; per gate \\
CI and tests pass & ci\_compliance\_checklist.md & IEC 62304; 172 tests; Python 3.10--3.12 \\
\multicolumn{3}{l}{\textit{Parent: physical-ai .../final\_paper/sections/}} \\
PSL and USL gates pass & psl\_usl\_standards.tex & Site PSL; robot USL minimum by procedure \\
\bottomrule
\end{tabularx}
\caption{Attestation clauses mapped to the recorded evidence and governing standards.}
\label{tab:attest-items}
\end{table}

No executable code, model weights, or training data sets are filed with the
attestation; they are referenced, versioned, and reproducible from a stated seed,
as Section~\ref{sec:supporting_documentation} specifies. The attestation files
the certification and the pointers; the records remain in the repository where a
reviewer can re-run them.
